\section{Provider Neutral, Capability Aware}

\migaki{} is provider-neutral, but it should not be provider-naive. Different providers expose different capabilities: prompt caching, explicit cache breakpoints, tool calling, remote MCP support, code execution, file search, web search, structured outputs, reasoning controls, batch APIs, rate-limit semantics, data-retention modes, region controls, and tracing integrations.

A lowest-common-denominator abstraction would erase useful execution capabilities. Instead, \mir{} should remain portable while lowering remains capability-aware:

\begin{lstlisting}[language={}]
type ProviderCapabilities = {
  supportsPromptCaching: boolean
  supportsExplicitCacheBreakpoints: boolean
  supportsToolCalling: boolean
  supportsRemoteMCP: boolean
  supportsStructuredOutputs: boolean
  supportsBatching: boolean
  supportsReasoningControls: boolean
  supportsZeroDataRetention?: boolean
  maxContextTokens?: number
  cacheTtlOptions?: string[]
}
\end{lstlisting}

Provider backends lower optimized \mir{} into concrete execution plans. In the near term, providers do not execute \mir{} directly. \migaki{} executes through adapters targeting provider APIs, gateways, workflow runtimes, or local inference backends.

Framework adapters are equally important. The OpenAI Agents SDK exposes the primitives \migaki{} needs for a first trajectory recorder: agents, model calls, tool calls, handoffs, tracing spans, and run configuration \citep{openai_agents_js,openai_agents_tracing}. The \texttt{migaki-openai-agents-js} package uses those primitives to record a reusable graph without forking the SDK deeply.

This is the desired adapter posture. \migaki{} should not require every agent application to be rewritten around \mir{} on day one. It should be able to wrap existing execution, observe it, and gradually introduce comparison, replay, and optimization where evidence supports those steps.
